%% notes-tex/multimodal/043-datasets-primer/sections/2-genotypes.tex

\section{Genotype classes and how each enters a model\secstatus{todo}}
\label{sec:genotypes}

The store holds four genotype classes, and the graph processors that turn a record into
model input treat them differently. What follows is the behavior of the code at the 043
build (\file{torchcell/data/graph_processor.py}), measured on the processed samples.

\notesec{Deletions}
Kemmeren, Sameith, Hu, Jackson, Nadal-Ribelles, Messner, Zelezniak and Ohya: every
record names one or two deleted genes (two for Sameith Dm). The genotype is the set of
systematic names, and the \texttt{Perturbation} processor marks those nodes in
\texttt{pert\_mask}. Jackson 2020 adds a construction identity (plate well and
barcode) below the gene name, so its 784 records hold 72 constructions of 12
genotypes.

\notesec{Isolates}
Caudal 2024 and the two Muenzner 2024 datasets: a natural isolate's genotype is its
variant calls against S288C, five perturbation types with the counts of
Table~\ref{tab:isolates}, between 195 and 7{,}025 perturbations per Caudal record.
The \texttt{Perturbation} processor marks every gene with any of them, so a Caudal
sample has 5{,}575 of 6{,}607 nodes in \texttt{pert\_mask}: the mask then separates
almost nothing, and a rule over perturbation types (for example loss of function and
absence only) is needed before isolates are trained on.

%% SOURCE: experiments/043-expression-proteome-morphology/scripts/primer_043_tables.py
%% -> tables/t3-isolates.tex
\begin{table}[htbp]
\centering
\scriptsize
\setlength{\tabcolsep}{3pt}
\caption[]{\textbf{Perturbations of an isolate record, by type.} Mean count per record
over the records of each dataset (\file{results/store_census.json}). Sequence variants
are named by the gene they fall in; natural loss of function, gene absence and gene
presence are released calls; the few deletions are engineered markers.}
\label{tab:isolates}
%% GENERATED by experiments/043-expression-proteome-morphology/scripts/primer_043_tables.py
%% SOURCE: results/store_census.json, results/cross_dataset_correlation.json,
%% results/within_dataset_ceilings.json and the 019 ceiling files named in
%% tables/provenance.json. Do not edit by hand.

\begin{tabular}{>{\raggedright\arraybackslash}p{24mm}r>{\raggedleft\arraybackslash}p{18mm}>{\raggedleft\arraybackslash}p{18mm}>{\raggedleft\arraybackslash}p{18mm}>{\raggedleft\arraybackslash}p{18mm}>{\raggedleft\arraybackslash}p{18mm}>{\raggedleft\arraybackslash}p{18mm}}
\toprule
Dataset & Records, n & Perturbations per record, min to max & Sequence variant, mean per record & Natural LoF, mean per record & Gene absence, mean per record & Gene presence, mean per record & Deletion, mean per record \\
\midrule
Caudal 2024 & 943 & 195 to 7{,}025 & 5{,}047 & 383 & 135 & 84 & 0 \\
Muenzner 2024 protein & 795 & 0 to 7{,}025 & 5{,}017 & 377 & 133 & 84 & 0 \\
Muenzner 2024 turnover & 55 & 4{,}903 to 6{,}430 & 5{,}059 & 373 & 134 & 90 & 0 \\
\bottomrule
\end{tabular}

\end{table}

\notesec{Segregants}
Boocock 2025 and N'Guessan 2025: a \texttt{SegregantGenotype} names a cross, two parent
strains pinned to sha256-verified assemblies, and a list of haplotype blocks, and it
carries no perturbation list by design. Since the stage-2 fix the processors read a
missing list as zero perturbed genes, so a segregant sample today has
\texttt{pert\_mask} all false: a model sees it as the reference strain. Using
segregants needs a representation of the blocks; Sec.~\ref{sec:proposal} names one.

\notesec{Promoter library}
Vaishnav 2022: each of 71{,}158 records is one 80 bp variable promoter insert driving a
YFP reporter, read as one scalar. Every record names the same reporter, so the store
holds one gene-name genotype for all of them; the perturbation sets separate them only
by sequence content. No gene-level vector and no gene-set encoder applies.
